\section{제안 방법}
\label{sec:method}

제안 절차는 네 단계로 구성된다. 첫째, 자료 정리 단계에서 같은 장면 식별 번호(자료 파일의 \texttt{clip\_id})를 가진 CoC를 시간순으로 놓고 각 값의 출처를 보존한다. 둘째, 규칙 변환 단계에서 문장 속 행동, 계속 지켜야 할 요구, 그 요구가 끝나는 조건과 행동 순서를 공통 형식으로 만든다. 셋째, 이전 요구를 기억하는 검사기가 사건마다 남은 요구를 갱신하고 오류 종류와 첫 위반 사건 번호를 기록한다. 넷째, 자연 자료, 합성 규칙 시험과 Python--UPPAAL 일치 결과를 서로 나누어 보고한다.

자료 정리 단계에서는 문장이 맞는지 또는 차량이 위험했는지를 판단하지 않는다. 센서에서 직접 얻은 입력, 기계가 매긴 품질 점수와 자차 속도에서 계산한 반응을 구분해 저장하고, 장면과 사건 순서를 고정한다. 규칙 변환 단계에서는 ``정지한다'', ``통과할 때까지'', ``그 뒤 진행한다''처럼 여러 단계가 있는 문장을 마지막 행동 하나로 줄이지 않는다. 대신 요구 행동, 계속되는 요구, 종료 조건과 선행 행동으로 나누고, 각 항목이 원문의 어디에서 왔는지와 문장을 해석하지 못한 이유를 남긴다.

검사기는 새 요구가 나오면 아직 지켜야 할 목록에 넣고, 요구가 끝났다는 명시적 증거가 참이면 목록에서 지운다. 종료 조건이 거짓이면 요구를 유지하고, 판단할 증거가 없으면 {\unknown}과 그 이유를 함께 남긴다. 마지막 단계에서는 자연 자료 판정, 합성 시험 사례 판정, Python--UPPAAL 일치, 보조 반응과 자료 출처 결과를 별도로 보고한다. 이렇게 하면 대상 수와 질문이 다른 결과가 하나의 성능 수치로 잘못 합쳐지는 것을 막을 수 있다.

\begin{figure*}[t]
\centering
\resizebox{0.84\textwidth}{!}{%
\begin{tikzpicture}[
  box/.style={draw, rounded corners=2pt, align=center, minimum height=11mm,
    text width=2.55cm, font=\footnotesize},
  small/.style={draw, rounded corners=2pt, align=center, minimum height=8mm,
    text width=3.15cm, font=\scriptsize},
  obs/.style={box, fill=observed!10},
  der/.style={box, fill=derived!10},
  ass/.style={box, fill=assumed!10},
  regression/.style={box, fill=neutral!8},
  route/.style={box, fill=neutral!8},
  note/.style={small, fill=black!3},
  flow/.style={-{Latex[length=2mm]}, thick, draw=black!80},
  feedback/.style={-{Latex[length=1.8mm]}, dashed, thick, draw=black!55}
]

% Four-component evidence architecture.
\node[obs] (collector) at (0,0) {Input organizer\\ordered CoC text\\times and data sources};
\node[ass] (compiler) at (4.0,0) {Checking-rule builder\\requirement, end condition\\action order, missing data};
\node[der] (events) at (7.4,0) {Fixed checking inputs\\shared format, XML\\checks};
\node[der] (observer) at (11.1,0) {Checker retaining\\prior requirements\\end condition, violation};
\node[route] (router) at (14.8,0) {Separated reports\\natural review\\auxiliary analyses};

\node[ass] (mut) at (7.4,-2.0) {Synthetic test pairs\\one changed item\\expected result required};
\node[der] (mutobserver) at (11.1,-2.0) {Current-event/prior-requirement\\checks on the same input};
\node[regression] (regression) at (14.8,-2.0) {Predefined rules exercised\\not detection accuracy};
\node[note] (boundary) at (0,-2.0) {No new driving scenes\\or control actions};

\node[route] (select) at (11.1,-4.0) {Selected natural sample\\independent review\\discussion afterward};
\node[note] (scope) at (7.4,-4.0)
  {Separate recorded, calculated, researcher-set, and modified evidence.};

\draw[flow] (collector) -- (compiler);
\draw[flow] (compiler) -- (events);
\draw[flow] (events) -- (observer);
\draw[flow] (observer) -- (router);
\draw[feedback] (events) -- (mut);
\draw[flow] (mut) -- (mutobserver);
\draw[flow] (mutobserver) -- (regression);
\draw[flow] (router.east) -- ++(1.8,0) |- (select.east);
\draw[flow] (select) -- (scope);

% Feedback is routed outside the node area to avoid crossing boxes.
\coordinate (feedbackleft) at ([xshift=-6mm]boundary.west);
\draw[feedback] (select.south) -- ++(0,-0.6) coordinate (feedbackbottom)
  -- node[midway,below=1mm,font=\scriptsize,align=center]{Reviewer feedback}
  (feedbackbottom -| feedbackleft) -- (feedbackleft) |- (collector.west);

\end{tikzpicture}%
}

\bifigcaption{자료 정리, 검사 규칙 변환, 이전 요구를 기억하는 판정과 결과 분리 절차. 자연 자료 검토와 기대 판정을 미리 확인한 합성 규칙 시험은 별도 경로로 유지한다.}{Data organization, rule construction, history-aware checking, and separate paths for natural review and controlled synthetic tests.}
\label{fig:method-pipeline}
\end{figure*}

그림~\ref{fig:method-pipeline}의 위쪽은 공통 입력을 검사 규칙과 판정 기록으로 바꾸는 과정이고, 아래쪽은 서로 다른 평가 자료를 뜻한다. 자연 자료에 대한 사람 판정은 합성 변형을 만들거나 검사 기준을 조정하는 데 사용하지 않는다. 반대로 합성 사례의 기대 판정도 자연 자료의 정답으로 사용하지 않는다. 따라서 40쌍이 정해 둔 규칙을 모두 실행했다는 결과와 자연 오류를 얼마나 잘 찾는지는 구분된다.

사건 $e_i$를 읽기 전에 아직 지켜야 하는 요구의 집합을 $A_{i-1}$이라 하면, 상태를 갱신하는 함수 $T$는 다음 요구 집합과 판정을 계산한다.
\begin{equation}
(A_i,V_i)=T(A_{i-1},e_i,q_i),
\label{eq:state-transition}
\end{equation}
여기서 $q_i\in\{\textsc{True},\textsc{False},\textsc{Unknown}\}$는 요구 종료나 외부 객체 상태에 대한 참ㆍ거짓ㆍ미상 증거이고, $V_i$는 정상 판정 또는 일관성 오류의 종류이다. 현재 사건만 보는 검사는 매번 이전 요구를 비우지만($A_{i-1}=\emptyset$), 이전 요구를 기억하는 검사는 이를 다음 사건으로 넘긴다. 이 차이 때문에 후자는 정지 요구가 끝나기 전에 진행하는 경우와, 끝난 정지 요구를 검사 목록에서 지우지 않은 경우를 판정할 수 있다.

\begin{table}[t]
\centering
\bitablecaption{이전 요구를 기억하는 검사가 확인하는 네 일관성 오류.}{Four consistency-error types checked while retaining prior requirements.}
\label{tab:contradiction-semantics}
\footnotesize
\begin{tabularx}{\columnwidth}{L{0.44\columnwidth}Y}
\toprule
Consistency error & Condition that causes it \\
\midrule
Hold--go conflict & Proceed while a stop or yield requirement remains \\
Requirement ended too early & End a requirement before its end condition is met \\
Action order violation & Start the next action before the required earlier action \\
Requirement left after its end & Keep a requirement after its end condition is met \\
\bottomrule
\end{tabularx}
\end{table}

같은 사건 목록은 \uppaal 검사 모델로도 변환된다. 이 모델은 사건을 하나씩 읽으면서 요구가 시작되고, 이어지고, 끝나는 과정과 위반 여부를 갱신한다. 한 사건의 상태 갱신이 끝날 때까지 다음 사건을 처리하지 않는 \uppaal{} 기능을 사용하여 사건 순서를 고정한다. 각 단계가 오류 없이 이어지는지와 오류 상태에 도달하는지를 확인하고, 처음 위반한 사건 번호를 실행 기록에 남긴다. 이 모델은 기록된 문장 순서만 검사하며 실제 차량 움직임이나 환경을 새로 만들어 시험하지 않는다.

Python 검사기와 \uppaal 검사 모델은 같은 사건과 증거를 입력받는다. Python은 장면별 판정과 남아 있는 요구의 변화를 기록하고, 변환 프로그램은 같은 내용을 UPPAAL 모델 파일에 넣는다. 두 결과를 비교할 때는 최종 정상/오류 판정, 오류 종류, 첫 위반 사건과 {\unknown} 사유를 확인한다. 예를 들어 정지 요구가 끝나기 전에 진행한 사례에서는 정지 요구가 시작되어 목록에 남고, 이후 진행 사건을 읽을 때 오류로 바뀌는 순서를 확인할 수 있다. 따라서 정상/오류만 표시할 때보다 어떤 사건에서 왜 오류가 났는지 추적하기 쉽다.

두 \uppaal 모델의 역할은 다르다. 중심 검사 모델은 연속 CoC의 순서와 남아 있는 요구를 확인하고, 그림~\ref{fig:uppaal-network}의 보조 모델은 \texttt{scene-0001}에서 반응 시간, 반복 사건과 자료 출처 처리를 확인한다. 표~\ref{tab:model-families}처럼 두 모델을 나누어 보고하여, 반응 시간 시험이 중심 문장 검사를 다시 검증한 결과로 오해되지 않게 한다.

\begin{table}[t]
\centering
\bitablecaption{평가에 사용한 두 UPPAAL 검사 모델과 각 모델의 확인 대상.}{Two UPPAAL checking models and what each checks.}
\label{tab:model-families}
\footnotesize
\begin{tabularx}{\columnwidth}{L{0.32\columnwidth}L{0.30\columnwidth}Y}
\toprule
Model & Checking components & What is checked \\
\midrule
Sequential CoC test cases & Requirement-memory checker & Order, error type/location, missing information \\
\texttt{scene-0001} response model & Event replay and three checks & Start/end, deadline, repeat clock, overwrite \\
\bottomrule
\end{tabularx}
\end{table}

이제 보조 반응 검사를 여러 역할로 나눈 \uppaal 모델을 설명한다.

\begin{figure*}[t]
\centering
\resizebox{0.78\textwidth}{!}{%
\begin{tikzpicture}[
  node distance=13mm and 15mm,
  box/.style={draw, rounded corners=2pt, align=center, minimum height=13mm,
    text width=3.25cm, font=\footnotesize},
  source/.style={box, fill=observed!12},
  monitor/.style={box, fill=derived!11},
  compiler/.style={box, fill=assumed!13},
  bus/.style={draw, rounded corners=2pt, fill=black!4, align=center,
    minimum width=11.7cm, minimum height=11mm, font=\footnotesize},
  flow/.style={-{Latex[length=2.3mm]}, line width=0.7pt, draw=black!75},
  event/.style={-{Latex[length=2.3mm]}, line width=0.85pt, draw=derived!80},
  audit/.style={-{Latex[length=2.3mm]}, line width=0.85pt, draw=assumed!85},
  label/.style={font=\scriptsize, align=center, fill=white, inner sep=1.5pt}
]

% Sources.
\node[source] (coc) at (0,3.4) {
  \textbf{CoC event replay}\\
  replay: CoC from stored data\\
  action-tagged decision events\\
  accepted / repeated / rejected
};

\node[source] (ego) at (7.7,3.4) {
  \textbf{Ego-motion replay}\\
  replay: ego motion from stored data\\
  selected speed-change condition\\
  calculated response evidence
};

% Researcher-selected contract semantics, kept distinct from recorded evidence.
\node[compiler] (compiler) at (3.85,5.45) {
  \textbf{Checking-rule builder}\\
  selected intent, $D$, repeat, $R_a$\\
  researcher-set checking rules
};

% Broadcast bus.
\node[bus] (bus) at (3.85,1.45) {
  \textbf{Fixed-order event replay}\\
  ordered decision events + calculated ego-response evidence\\
  requirement start, repeat, response, rejected input
};

% Monitors.
\node[monitor] (obl) at (-1.6,-0.95) {
  \textbf{Requirements across events}\\
  start / retain / end\\
  Idle $\rightarrow$ Active\\
  Active $\rightarrow$ Responded
};

\node[monitor] (react) at (3.85,-0.95) {
  \textbf{Response timing}\\
  clock $reaction$\\
  deadline $D$\\
  repeated-event rule
};

\node[monitor] (valid) at (9.3,-0.95) {
  \textbf{Input quality}\\
  quality check\\
  overwrite check
};

% Compiler selects contract semantics; sources replay the resulting evidence.
\draw[audit] (compiler.south west) -- node[label, left] {selected\\rule} (coc.north);
\draw[audit] (compiler.south east) -- node[label, right] {response\\condition} (ego.north);

% Source to bus.
\draw[event] (coc.south) -- node[label, left] {ordered decision\\events} (bus.north west);
\draw[event] (ego.south) -- node[label, right] {calculated ego\\response evidence} (bus.north east);

% Bus to monitors. Route from distinct bottom anchors to avoid visual crossing.
\draw[event] ([xshift=-4.2cm]bus.south) -- node[label, left] {start\\response} (obl.north);
\draw[event] (bus.south) -- node[label, right] {start\\repeat\\response} (react.north);
\draw[event] ([xshift=4.2cm]bus.south) -- node[label, right] {rejected input} (valid.north);

% Result box.
\node[draw, rounded corners=2pt, fill=neutral!8, align=center,
  minimum width=11.7cm, minimum height=10mm, font=\footnotesize] (queries)
  at (3.85,-3.05) {
  \textbf{Results}: rule pass / review / data source / modified case\\
  states checked: requirement active, response seen, deadline missed,
  repeat-clock reset, rejected-data overwrite, stopped execution
};

\draw[flow] (obl.south) -- (queries.north west);
\draw[flow] (react.south) -- (queries.north);
\draw[flow] (valid.south) -- (queries.north east);

\end{tikzpicture}%
}

\bifigcaption{\texttt{scene-0001}의 기록 사건 재생과 요구ㆍ반응ㆍ자료 출처 검사를 역할별로 나눈 UPPAAL 모델.}{UPPAAL model separating fixed event replay from requirement, response, and data-source checks for scene-0001.}
\label{fig:uppaal-network}
\end{figure*}

그림~\ref{fig:uppaal-network}에서 CoC 사건 재생 부분은 최초 판단, 반복 판단과 품질 기준을 통과하지 못한 사건을 시간순으로 읽고, 자차 움직임 재생 부분은 속도에서 찾은 반응을 읽는다. 세 검사 부분은 각각 요구의 시작ㆍ유지ㆍ종료, 반응 제한 시간과 반복 시계, 낮은 품질 사건이 기존 정보를 덮어쓰는지를 확인한다. 이 모델은 새 행동이나 환경을 만들지 않고, 기록된 사건을 읽었을 때 검사 상태가 의도한 순서로 바뀌는지만 확인한다.

표~\ref{tab:scene1-map}은 원래 기록된 시각과 검사 모델이 읽는 사건의 관계를 보여 준다. 첫 감속 CoC에서 지켜야 할 요구가 시작되고, 두 번의 반복 CoC는 같은 요구에 연결된다. 자차 속도에서 계산한 반응이 확인되면 그 요구를 끝낸다. 마지막의 낮은 품질 사건은 앞서 받아들인 정보를 덮어쓰지 않아야 한다. 센서가 기록한 값, 기록에서 계산한 값, 연구자가 정한 값을 표에서 구분하여 각 값의 출처를 확인할 수 있게 했다.

\begin{table}[t]
\centering
\bitablecaption{\texttt{scene-0001}의 기록값ㆍ계산값과 UPPAAL 사건의 대응 관계.}{Mapping recorded and calculated scene-0001 evidence to UPPAAL events.}
\label{tab:scene1-map}
\scriptsize
\begin{tabularx}{\columnwidth}{L{0.34\columnwidth}L{0.25\columnwidth}Y}
\toprule
Recorded or calculated evidence & Event given to the model & Meaning \\
\midrule
Accepted deceleration at 2.500 s & First requirement at 0.000 s & Start the requirement \\
Repeated deceleration at 3.500 s & Repeated requirement at 1.000 s & Record a repeated event \\
Speed-decrease onset at 4.660 s & Response at 2.160 s & End with the calculated response \\
Repeated deceleration at 4.800 s & Repeated requirement at 2.300 s & Record a repeat after response \\
Low-quality CoC at 9.300 s & Rejected input at 6.800 s & Do not overwrite accepted data \\
\bottomrule
\end{tabularx}
\end{table}

자연 자료와 별도로 네 일관성 오류마다 정상 기준 사례 10개를 만들고, 각 기준 사례에서 한 항목만 바꾼 사례를 짝지었다. 생성 프로그램은 기준 사례에는 오류가 없고 바꾼 사례에는 미리 지정한 오류 하나만 나타나는지 확인한 뒤에만 그 쌍을 사용했다. 따라서 합성 40쌍은 정해 둔 네 규칙이 의도대로 작동하는지 확인하는 시험 사례이다. 기대 판정이 나오도록 생성 단계에서 미리 걸렀으므로, 처음 보는 자연 오류를 얼마나 잘 찾는지 평가하는 자료는 아니다.

한 항목만 바꾸면 어떤 변화가 어떤 오류를 만들었는지 분명해진다. 첫째, 정지ㆍ양보 요구가 끝났다는 증거 없이 진행을 시작한다. 둘째, 종료 조건이 아직 거짓인데 요구를 끝낸다. 셋째, 꼭 먼저 해야 하는 행동보다 다음 행동을 앞세운다. 넷째, 종료 조건이 참인데도 끝난 요구를 목록에서 지우지 않는다. 두 검사에는 같은 현재 사건 정보를 주고, 현재 사건만 보는 검사에만 이전 요구를 전달하지 않았다. 따라서 판정 차이는 문장 해석이나 입력 차이가 아니라 이전 요구를 기억하는지에서 생긴다.

Python과 UPPAAL 비교에는 정상 사례 하나, 네 오류 사례, 증거 부족 사례와 문장 해석 실패 사례 등 일곱 표준 시험 사례를 사용한다. 최종 판정, 오류 종류, 첫 위반 사건과 {\unknown} 사유를 비교한다. 두 구현이 같은 판정 규칙을 사용하므로, 이 비교는 규칙 자체가 옳다는 독립 증명보다 두 구현이 같은 결과와 기록을 내는지 확인하는 시험이다.

\begin{table}[t]
\centering
\bitablecaption{Python--UPPAAL 구현 일치 확인에 사용한 표준 시험 사례.}{Standard test cases for Python--UPPAAL implementation agreement.}
\label{tab:uppaal-fixtures}
\footnotesize
\begin{tabularx}{\columnwidth}{L{0.28\columnwidth}cY}
\toprule
Test-case group & Count & Compared output \\
\midrule
Consistent sequence & 1 & Consistent decision \\
Four consistency errors & 4 & Error type and first violating event \\
Missing evidence/text interpretation & 2 & Insufficient-information decision and reason \\
\midrule
Total & 7 & Final decision and decision-record items \\
\bottomrule
\end{tabularx}
\end{table}
